\section*{はじめに}
\label{sec:introduction}
  \subsection*{研究の背景と目的}
    コンピュータは四則演算を行うことを前提として設計されているため，多項式の計算は高速に行うことができる．
    しかし，多項式で直接あらわすことができない関数（例えば指数関数や三角関数など）というのは多々存在する．

    そのため，そういった関数を多項式で近似することで，計算機上での数値計算の実行速度を高速化できるなどのメリットが生まれる．
    さて，近似多項式の作成にはいくつか知られた手法が存在する．
    例えば，テイラー展開は，指定した点付近での精度は極めて高いが，その点から離れるにつれて誤差が増大する．そのため，定義域内においての最大誤差の抑制が求められる場合には不向きである
    その対策として事前に設定した標本点を通る多項式であるラグランジュ補間多項式は，閉区間 $\left[a,b\right]$ における$n$次多項式を一意に決定することができる．
    しかし，煩雑な標本点を用いた場合など標本点の取り方によって求める誤差関数の最大ノルムよりも大きな誤差が生じる場合がある．
    そのため，ここではラグランジュ補間多項式を用いた多項式近似において誤差関数の最大ノルムを最小化した最良近似多項式とそれを取得するために必要となる標本点の導出手法，それによってどのような近似を得ることができるかについて考える．

  \subsection*{本論文の構成}
    \begin{list}{}{
        \setlength{\leftmargin}{4.5em} % 本文のインデント幅（「第X章」の幅+α）
        \setlength{\labelwidth}{4.5em} % ラベル「第X章」の幅
        \setlength{\itemindent}{0pt}   % 1行目のインデント補正
        \setlength{\labelsep}{0.5em}   % ラベルと本文の間隔
        \setlength{\itemsep}{1em}      % 項目ごとの隙間（ここを調整して広げられます）
    }

    \item[\textbf{第1章}]
    「ラグランジュ補間多項式とその評価」では，関数近似の基礎となる多項式補間の理論について述べる．
    特にラグランジュ補間について解説し，等間隔な標本点を用いた場合に発生するルンゲ現象の問題点について，具体的な数値例を用いて示す．

    \item[\textbf{第2章}]
    「最良近似多項式の理論」では，本研究の核心である最良近似の数学的基礎について述べる．
    一様ノルムによる誤差の定義，最良近似多項式の存在と一意性，そしてRemesアルゴリズムの収束判定条件となるチェビシェフの交替定理について解説する．

    \item[\textbf{第3章}]
    「Remesアルゴリズムによる最良近似多項式の計算」では，最良近似多項式を数値的に導出する手法について述べる．
    特に，反復ごとに参照点を更新する「全点同時入れ替え法」の手順と，その2次収束性について理論的な背景を説明する．

    \item[\textbf{第4章}]
    「極値点探索の手法」では，Remesアルゴリズムの収束性能を左右する誤差関数の極値探索に焦点を当てる．
    微分を用いずに高速な収束を実現するBrent法のアルゴリズム構成と，その根拠について説明する．

    \item[\textbf{第5章}]
    「Brent法を用いたRemesアルゴリズムの実装と評価」では，提案手法の実装詳細と数値実験の結果について述べる．
    ルンゲ関数 $f(x) = 1/(1+25x^2)$ を対象とした近似実験を行い，グリッド極値点探索を用いた従来手法と比較して，提案手法が機械イプシロン程度の精度で交替定理を満たす解を導出できることを定量的に示す．

    \item[\textbf{第6章}]
    「考察」では，数値実験の結果に基づき，グリッド極値点探索法とBrent法の性能差が生じた要因について数理的な観点から議論する．
    特に，各極値における残留誤差の定量評価を通じ，Brent法の導入がRemesアルゴリズムの理論的な2次収束性を実現するために不可欠であることを明らかにする．

    \item[\textbf{第7章}]
    「結論」では，本研究で得られた知見を総括し，今後の課題と展望について述べる．
    \end{list}
  \subsection*{謝辞}
    本研究を進めるにあたり，多大なるご指導とご助言をいただきました指導教員の先生に深く感謝いたします．
    また，有益な助言をいただいた研究室の皆様，そして在学中支えてくれた家族に心より感謝申し上げます．
\newpage  %% 改ページ ==================================================